\section{Introduction}
\label{sec:introduction}

Reaction--diffusion systems with gradient-dependent nonlinearities arise when local production, dissipation, or transfer mechanisms depend on spatial variations of the state. We consider
\begin{equation}
\begin{cases}
\partial_tu_i-d_i\Delta u_i=f_i(t,x,u,\nabla u), &(t,x)\in Q_T=(0,T)\times\Omega,\\
u_i(0,x)=u_{i,0}(x),&x\in\Omega,\\
u_i(t,x)=0,&(t,x)\in(0,T)\times\partial\Omega,
\end{cases}
\qquad i=1,\ldots,m,
\label{eq:main_system}
\end{equation}
where $\Omega\subset\mathbb R^N$ is bounded and $d_i>0$. Quadratic dependence on $\nabla u$ is critical for the compactness and equi-integrability arguments used to construct weak solutions. At the modelling level, such terms describe rates that become stronger near steep spatial fronts and can therefore serve as reduced representations of interfacial transfer, gradient-enhanced dissipation, or front-localized conversion. The present work studies a mathematically representative prototype; application-specific constitutive validity would require dimensional derivation, parameter calibration, and observational validation.

\subsection{Related work and research gap}
\label{subsec:related-work}

\paragraph{Analytical background.}
Global existence and boundedness for reaction--diffusion systems under mass control or mass dissipation have been developed through several complementary mechanisms. Classical results established global existence under structural assumptions linking reaction terms and diffusion~\cite{hollis1987global,boudiba2000coupled}, while later work clarified the role of mass control, quadratic growth, duality, and intermediate-sum structures~\cite{pierre2010survey,souplet2018quadratic,fellner2020quadratic,cupps2021uniform,morgan2020intermediate}. Critical natural growth with respect to the gradient is substantially more delicate because weak compactness of the state variables alone does not automatically control nonlinear expressions involving $|\nabla u|^2$. Scalar critical-gradient problems were treated by Landes and Mustonen~\cite{landes1994parabolic}; for systems, Alaa and Mounir considered a two-component critical-gradient setting~\cite{alaa2001global}, and triangular $m$-component extensions with low-regularity data were subsequently developed in~\cite{alaa2013triangular,bouarifi2015climate}.

The published scope of the two triangular references is important for positioning the present result accurately. Reference~\cite{alaa2013triangular} establishes weak-solution existence for $m\times m$ triangular systems with nonnegative solutions, $L^1$ data, and critical gradient growth, whereas~\cite{bouarifi2015climate} develops a related critical-gradient triangular existence framework in an energy-balance climate setting. Those references provide the theorem-level context. Here we isolate one local coercivity step, re-derive it from the regularized equations, and state exactly which full-gradient information that step requires. The resulting contribution is specific: the exact cross-gradient identity, the exposed full-gradient hypothesis, and the conditional removal of diffusion ordering at the absorption stage.

The analytical gap is a precise one inside the existing global-existence framework. When the first $r$ regularized equations are summed and tested by $T_k(U_{r,n})$, direct integration by parts produces cross terms containing the \emph{full} gradients $\nabla u_{j,n}$ of preceding components. Replacing them by $\nabla T_k(u_{j,n})$ changes the localization set and therefore changes the estimate being proved. The resulting question is precise: once full $L^2(Q_T)$ bounds for the preceding gradients are independently available, is an ordering of the positive diffusion coefficients still needed to absorb the cross terms? We show that it is not: conditional on the independently available full-gradient bounds, a free Young parameter absorbs all such terms for any fixed positive diffusion vector. The novelty does not lie in Young's inequality itself. It lies in identifying the exact cross-gradient structure generated by the partial-sum test, isolating the full-gradient hypothesis that is genuinely required, and showing that once this hypothesis is available no monotone ordering of the diffusion coefficients is needed at the coercivity stage. The contribution is therefore the exact identity together with this conditional coercivity clarification and an explicit separation from the additional compactness, equi-integrability, and nonlinear limit-identification arguments required for existence.

\paragraph{Numerical verification and classical discretizations.}
Because the analytical claim is paired with numerical evidence, implementation verification must be separated from comparison against a refined numerical reference. For parabolic problems, standard Galerkin theory and time-discretization analysis provide the classical context for $P_1$ finite elements and implicit--explicit reaction--diffusion updates~\cite{thomee2006galerkin,hundsdorfer2003adr}. We therefore use the Method of Manufactured Solutions (MMS), a standard code-verification methodology in which a smooth analytical field is prescribed and the corresponding forcing terms are obtained by substitution into the governing equations~\cite{roache2002mms,oberkampf2010vv}. Separate spatial and temporal refinements then test whether the implemented FDM and FEM recover their expected asymptotic behavior, independently of the unforced reference computation used later for method comparison.

\paragraph{Physics-informed neural networks.}
Physics-informed neural networks (PINNs) embed differential-equation residuals and data constraints into the training objective~\cite{raissi2019physics}; software frameworks and broader scientific-machine-learning surveys have accelerated their use for forward and inverse differential problems~\cite{lu2021deepxde,cuomo2022sciml}. Their flexibility, however, does not make optimization or error control automatic. Training can suffer from severe loss-gradient imbalance and unequal convergence of objective components~\cite{wang2021gradient}, neural-tangent-kernel analyses reveal mechanisms behind training failure~\cite{wang2022ntk}, and empirical studies show that a small training loss may coexist with a substantial solution error~\cite{krishnapriyan2021failure}. Rigorous analyses available for selected PDE classes likewise distinguish training, generalization, and field errors and require stability, regularity, and sampling assumptions to connect them~\cite{deryck2022error,mishra2023generalization,deryck2024numerical}. For time-dependent PDEs, causality-aware strategies can improve training~\cite{wang2024causality}, while residual-based adaptive sampling has been shown to materially affect accuracy and sampling efficiency~\cite{wu2023sampling}. Exact imposition of Dirichlet conditions through the network output is also an active design choice rather than a cosmetic one~\cite{berrone2023dirichlet}.

These observations define the computational research gap of the present paper. A credible comparison should not use a PINN training loss as its only validation quantity, nor should it compare an unverified classical implementation against a learned surrogate. We therefore combine: (i) independent MMS verification of FDM and FEM; (ii) a common refined-reference comparison for PINN, FDM, and FEM; (iii) a common-point residual test excluded from PINN training; (iv) positivity, partial-mass, and energy diagnostics; and (v) a three-component non-monotone benchmark that directly exercises the multiple cross terms occurring in the analytical estimate.

\subsection{Contributions and scientific architecture}

The contributions are:
\begin{enumerate}
 \item the exact partial-sum identity and a parametric Young estimate valid for any fixed positive diffusion vector under explicit preceding-gradient bounds;
 \item an explicit statement of the logical dependence between the exact estimate and the compactness theory developed in earlier work;
 \item an independent manufactured-solution verification showing approximately second-order spatial and first-order temporal behaviour for the classical discretizations;
 \item a reproducible two-component PINN--FDM--FEM diagnostic comparison with five-seed variability, independent residual testing, explicit initial-data error, refinement checks, mass and energy diagnostics, and a collocation-size ablation;
 \item a three-component non-monotone benchmark with two simultaneous cross terms and energy-derived preceding-gradient bounds;
 \item a constraint-compatible PINN representation that enforces non-negativity and homogeneous Dirichlet values exactly while leaving the initial condition penalized and quantitatively auditable.
\end{enumerate}

\begin{figure}[H]
\centering
\begin{tikzpicture}[
  scale=.82,
  transform shape,
  node distance=7mm and 6mm,
  box/.style={draw, rounded corners=2pt, align=center, minimum height=10mm, text width=2.75cm, inner sep=4pt, font=\small},
  emph/.style={box, line width=0.9pt},
  arrow/.style={-{Latex[length=2mm]}, line width=0.7pt},
  dashedarrow/.style={-{Latex[length=2mm]}, dashed, line width=0.7pt}
]
\node[box] (structure) {Triangular reaction--diffusion structure\\critical gradient growth};
\node[emph, right=of structure] (identity) {Exact partial-sum identity\\full preceding gradients};
\node[emph, right=of identity] (coercivity) {Conditional coercivity\\free Young parameter\\fixed $d_i>0$};

\node[box, below=of structure] (mms) {MMS code verification\\FDM + $P_1$ FEM\\space/time refinement};
\node[box, right=of mms] (two) {$m=2$ common benchmark\\PINN--FDM--FEM\\independent residual};
\node[emph, right=of two] (three) {$m=3$ non-monotone benchmark\\two simultaneous cross terms\\energy-based $C_1,C_2$};

\draw[arrow] (structure) -- (identity);
\draw[arrow] (identity) -- (coercivity);
\draw[arrow] (mms) -- (two);
\draw[arrow] (two) -- (three);
\draw[dashedarrow] (coercivity) -- (three);
\draw[dashedarrow] (structure) -- (mms);
\end{tikzpicture}
\caption{Scientific architecture of the paper.}
\label{fig:scientific-architecture}
\end{figure}

Figure~\ref{fig:scientific-architecture} summarizes the logic of the paper. The upper path represents analytical implication, whereas the lower path represents numerical verification and diagnostic comparison. The two paths are deliberately kept logically distinct and meet only at the level of a controlled numerical illustration: the three-component benchmark realizes the non-monotone diffusion configuration and independently available preceding-gradient estimates without being used as a proof of the analytical proposition.

Section~\ref{sec:theoretical} gives the analytical derivation. Section~\ref{sec:numerical} presents the common numerical protocol, Section~\ref{sec:results} reports the computed results and ablations, and Section~\ref{sec:conclusion} concludes.
